\chapter{Solid State Theory: Midterm II, May 11, 2026}
\label{q:theory-2026-05-midterm-ii}
\begin{sourceNote}
All seven questions are reconstructed from a six-page document containing both questions and supplied answers. The printed weights total 100\%. Answers have been separated into the companion solution book. The original PDF is not reproduced. The source's terminology and formulas are retained; qualifications concerning the ionic gap and the zero-frequency dielectric pole are discussed in the solutions.
\end{sourceNote}

\section*{1. Evjen estimate (10\%)}
Estimate the Madelung constant $\alpha_M$ of the NaCl structure by the Evjen method, considering an Evjen cell that is a cube of edge $a$, where $a$ is the lattice constant. You are given
\[
\alpha_M\approx\frac{n_1}{1}+\frac{n_2}{\sqrt2}
+\frac{n_3}{\sqrt3}+\frac{n_4}{\sqrt4}.
\]
What are the integers $n_1,n_2,n_3,n_4$?

\section*{2. Ionic radii and crystal structure (10\%)}
Consider an ionic crystal composed of two kinds of ions with radii $R_A$ and $R_C$.
\begin{enumerate}[label=(\alph*)]
\item What crystal structure does the crystal tend to adopt if $R_A=R_C$? Why?
\item Is it possible for the crystal to have the CsCl structure if
$R_A/R_C>(\sqrt3+1)/2$? Justify your answer.
\end{enumerate}

\section*{3. Ionic crystals (25\%)}
\begin{enumerate}[label=(\alph*)]
\item (5\%) Express the semi-empirical form of the cohesive energy per molecule $u(R)$, the specific cohesive energy $E_{\rm coh}$, of ionic crystals in terms of parameters $\lambda$ and $n$.
\item (10\%) Determine the bulk modulus $B_0$ of NaCl from the experimentally obtained equilibrium nearest-neighbor distance $R_0$ and specific cohesive energy $E_{\rm coh}$. Express $B_0$ in terms of $R_0$ and $E_{\rm coh}$, and/or the Madelung energy $E_M=-\alpha_M e^2/R_0$, instead of $\lambda$ and $n$.
\item (10\%) Estimate the band gap $\varepsilon_G$ of NaCl from the electron affinity of the Cl atom $\Delta(\mathrm{Cl}^-)$, the ionization energy of the Na atom $\Delta(\mathrm{Na})$, the equilibrium nearest-neighbor distance $R_0$, and the Madelung constant $\alpha_M$.
\end{enumerate}

\section*{4. Exciton absorption peaks (10\%)}
In the low-temperature absorption spectrum of a pure semiconductor, the lowest two exciton peaks have energies $\Delta_1$ and $\Delta_2$, with $\Delta_1<\Delta_2$. Determine the band gap and the largest exciton binding energy.

\section*{5. Complete static screening (10\%)}
Show that screening is complete in metals where the static dielectric function tends to infinity as $q\to0$.

\section*{6. Thomas--Fermi model (10\%)}
\begin{enumerate}[label=(\alph*)]
\item Express the Thomas--Fermi dielectric function in terms of the Thomas--Fermi wave number $k_{\rm TF}$.
\item Show that the model predicts an infinite screening electron density at an external point charge.
\end{enumerate}

\section*{7. Dynamic dielectric function of a metal (25\%)}
\begin{enumerate}[label=(\alph*)]
\item (5\%) From the long-wavelength dynamic dielectric function
\[
\epsilon(q,\omega)=1-\frac{\omega_p^2}{\omega^2}
-\frac{3\omega_p^2}{5\omega^4}v_F^2q^2,
\]
find the plasmon dispersion relation.
\item (5\%) Neglecting dispersion but including damping, the dielectric function is
\[
\epsilon(\omega)=1-\frac{\omega_p^2}{\omega^2+\ii\gamma\omega}.
\]
Find the energy-loss function.
\item (5\%) Write down the Kramers--Kronig relations.
\item (10\%) Does the dielectric function in (b) obey the Kramers--Kronig relations? If so, justify your answer. If not, modify the dielectric function so that it does.
\end{enumerate}
